% =====================================================================
%  Plantilla de documentación tipo poster — Portafolio Personal
%  Adaptada del template beamerposter (gemini/msu) con paleta naranja
%  accent del portafolio. Single page A3 landscape, multicolumna,
%  blocks con header coloreado.
% =====================================================================

\documentclass[10pt]{article}

% ------------------ Tamaño de página: custom 42x37cm (A3 ancho + más alto)
%   para que entre todo el contenido en una sola hoja con multicols. ------
\usepackage[paperwidth=42cm, paperheight=37cm, margin=1cm, top=1.2cm, bottom=1.2cm]{geometry}

% ------------------ Codificación y fuentes ------------------
% Con xelatex usamos fontspec; DejaVu Sans es nativo en Linux y se ve limpio.
\usepackage{fontspec}
\setmainfont{DejaVu Sans}
\setsansfont{DejaVu Sans}
\setmonofont{DejaVu Sans Mono}
\renewcommand{\familydefault}{\sfdefault}

% ------------------ Color, gráficos, layout ------------------
\usepackage{xcolor}
\usepackage{graphicx}
\usepackage{multicol}
\usepackage[most]{tcolorbox}
\usepackage{enumitem}
\usepackage{microtype}
\usepackage{hyperref}

% ------------------ Paleta del portafolio ------------------
% Colores: los placeholders {{COLOR_*}} los reemplaza generate-docs.ts según el
% tema. Para light: bodytext oscuro sobre fondo claro. Para dark: bodytext claro
% sobre fondo oscuro. El accent naranja se mantiene en ambos.
\definecolor{accent}{HTML}{F97316}
\definecolor{accentdark}{HTML}{{{COLOR_ACCENT_DARK}}}
\definecolor{accentlight}{HTML}{{{COLOR_ACCENT_LIGHT}}}
\definecolor{bodytext}{HTML}{{{COLOR_BODY_TEXT}}}
\definecolor{mutedtext}{HTML}{{{COLOR_MUTED_TEXT}}}
\definecolor{cardbg}{HTML}{{{COLOR_CARD_BG}}}
\definecolor{pagebg}{HTML}{{{COLOR_PAGE_BG}}}
\pagecolor{pagebg}
\color{bodytext}

\hypersetup{colorlinks=true, urlcolor=accent, linkcolor=accent}

% ------------------ Bloque con header naranja ------------------
\newtcolorbox{posterblock}[1]{
  enhanced,
  colback=cardbg,
  colframe=accent,
  coltitle=white,
  coltext=bodytext,   % por defecto tcolorbox usa negro; forzar bodytext para dark mode
  colbacktitle=accent,
  title={\strut\large\bfseries\MakeUppercase{#1}},
  boxrule=0pt,
  toprule=0pt,
  bottomrule=0pt,
  leftrule=0pt,
  rightrule=0pt,
  arc=3pt,
  outer arc=3pt,
  left=9pt, right=9pt, top=6pt, bottom=6pt,
  toptitle=3pt, bottomtitle=3pt,
  before skip=0pt,
  after skip=10pt,
  fonttitle=\sffamily,
  breakable
}

% ------------------ Subtítulo dentro de block (heading) ------------------
\newcommand{\posterheading}[1]{%
  \vspace{4pt}%
  \par{\color{accentdark}\sffamily\bfseries #1}\par
  \vspace{2pt}%
}

% ------------------ Footer / pagestyle ------------------
\pagestyle{empty}

% ------------------ Espaciado entre items ------------------
\setlist{itemsep=2pt, topsep=2pt, leftmargin=1.2em}

\begin{document}
\sloppy

% ==================== ENCABEZADO ====================
% Los placeholders {{...}} los reemplaza backend/scripts/generate-docs.ts antes
% de compilar. Si compilas la plantilla a mano (sin el script), los placeholders
% aparecen literales en el PDF — usá el script para regenerar.
\noindent
\begin{minipage}{0.7\textwidth}
  \raggedright
  {\color{accent}\fontsize{42}{50}\selectfont\bfseries {{TITLE}}\par}
  \vspace{6pt}
  {\color{mutedtext}\Large {{SUBTITLE}}}
\end{minipage}
\hfill
\begin{minipage}{0.28\textwidth}
  \raggedleft\small\color{mutedtext}
  \textbf{Versión} {{VERSION}} \\
  \textbf{Última actualización} {{LAST_UPDATE}} \\
  \textbf{Repositorio} \href{{{REPO_URL}}}{{{REPO_LABEL}}} \\
  \textbf{Demo} \href{{{DEMO_URL}}}{{{DEMO_LABEL}}}
\end{minipage}

\vspace{6pt}
\noindent\textcolor{accent}{\rule{\linewidth}{2pt}}
\vspace{14pt}

% ==================== CONTENIDO EN 3 COLUMNAS ====================
\setlength{\columnsep}{18pt}
\begin{multicols}{3}

% ----- COLUMNA 1 -----

\begin{posterblock}{Resumen}
{{OVERVIEW}}

\posterheading{Características principales}
{{FEATURES_LIST}}
\end{posterblock}

\begin{posterblock}{Stack tecnológico}
\posterheading{Core}
React 18 sobre TypeScript 5.6 con modo estricto, empaquetado con Vite 5.4.

\posterheading{Estilos y animaciones}
Tailwind CSS, lottie-react para la animación del hero, CSS 3D transforms para las carpetas, simulación force-directed propia en SVG para el grafo.

\posterheading{APIs del navegador}
Web Audio API con GainNode, localStorage, Pointer Events para drag y pinch zoom, Page Visibility para pausar música, MediaQuery para detectar hover y preferencia de tema.

\posterheading{Infraestructura}
pnpm como gestor obligatorio, GitHub Actions para CI, GitHub Pages para hosting estático global con HTTPS.

\posterheading{Bundle final}
Aproximadamente 180 kilobytes comprimidos. Sin React Router, Framer Motion ni librerías de state management.
\end{posterblock}

% ----- COLUMNA 2 -----

\begin{posterblock}{Arquitectura del frontend}
La aplicación se estructura en cinco capas verticales que separan responsabilidades: herramientas de desarrollo, framework principal, estilos e interfaz, APIs del navegador y despliegue continuo. Cada capa solo depende de las superiores y expone una interfaz acotada a las inferiores.

\posterheading{Capas verticales}
Dev/Build (pnpm, TypeScript, Vite) que solo dependen de las herramientas del sistema. Core (React, Hash Routing) que monta el sitio. Estilos y UI (Tailwind, lucide, lottie, CSS 3D) que dan la identidad visual. Browser APIs (Web Audio, localStorage, Pointer Events, Page Visibility, MediaQuery) que aprovechan capacidades nativas sin libs adicionales. CI/Deploy (GitHub Actions y Pages) que lleva el bundle compilado a producción.

\posterheading{Decisión clave}
Hash routing custom de aproximadamente treinta líneas reemplaza a react-router, lo que elimina una dependencia significativa y mantiene URLs compartibles que sobreviven a refrescos del navegador. Ver el diagrama completo del stack en la sección Diagramas del portafolio.
\end{posterblock}

\begin{posterblock}{Componentes principales}
\posterheading{FolderPortfolio}
Renderiza el grid de seis carpetas tridimensionales con abanico de cards y orquesta el lightbox modal.

\posterheading{ProjectDetailView}
Vista de detalle con hero, grafo conceptual force-directed y seis tarjetas de sección. Navegación entre proyectos sin volver al grid.

\posterheading{ConceptGraph}
Simulación física propia con repulsión entre nodos, springs en aristas y levitación continua. Soporta drag con pointer events y filtro por grupo.

\posterheading{DiagramsViewer}
Visor fullscreen con zoom hasta ochocientos por ciento, pan y pinch zoom en mobile. Navegación circular entre múltiples imágenes.

\posterheading{PortfolioChat}
Burbuja flotante con panel desplegable. Mock por keyword matching, integración con LLM planeada como próxima fase.
\end{posterblock}

% ----- COLUMNA 3 -----

\begin{posterblock}{Arquitectura del backend}
Planificación de una capa serverless mínima que coexiste con el frontend estático sin convertirlo en monolito. Dos workers en runtime para chat y administración, una suite de agentes batch en GitHub Actions para auto-generar contenido del portafolio.

\posterheading{Tres capas funcionales}
Eventos (push a repo, release tag, cron, chat de visitante, aprobación de admin) que disparan la cadena. Backend (workers de runtime para chat y admin de borradores, agentes batch en GitHub Actions agrupados en content, lifecycle y distribution, persistencia mixta entre JSON committed y Workers KV ephemeral). Outputs (portafolio web auto-desplegado, CV PDF auto-generado, blog en Dev.to, post en LinkedIn, respuesta del chat).

\posterheading{Decisión clave}
JAMstack híbrida en lugar de monolito Express o n8n. El tráfico impredecible de un portafolio personal justifica pagar por request en Cloudflare Workers en lugar de mantener un servidor 24/7. Ver el diagrama completo del flujo en la sección Diagramas del portafolio.
\end{posterblock}

\begin{posterblock}{Modelo de datos}
El catálogo de proyectos vive en \texttt{src/data/projects.ts} como un arreglo de seis categorías. Cada proyecto declara identificador único, imagen, título y descripción bilingües, lista de tags, estado dentro de tres valores, enlaces a repositorio y demo, conceptos del grafo y categorías cross-disciplina opcionales.

\posterheading{Tipos clave}
\texttt{LocalizedString}, \texttt{ProjectStatus}, \texttt{ConceptGroup} con cuatro valores semánticos: arquitectura, datos, operaciones y machine learning.

\posterheading{Persistencia}
JSON y markdown committed al repositorio para datos públicos. Workers KV para rate limiting del chat y cola de borradores de distribución. No se almacena historial de conversaciones por privacidad.
\end{posterblock}

\begin{posterblock}{Próximos pasos y cambios recientes}
\posterheading{Cambios recientes}
{{RECENT_CHANGES}}

\posterheading{Próximos pasos}
{{ROADMAP_LIST}}
\end{posterblock}

\end{multicols}

\end{document}
